\documentclass[Serif, 10pt, brown]{beamer}
\usepackage{booktabs,xcolor}
\usepackage{pifont}
\newcommand*\CHECK{\ding{51}}
\usepackage{array}
\newcolumntype{P}[1]{>{\centering\arraybackslash}p{#1}}

\usepackage{setspace,mathtools,amssymb,multirow,array,amsmath,tikz}
\usepackage[normalsize]{subfigure}
\usetikzlibrary{patterns,automata,positioning,decorations.pathreplacing,decorations,calc}
\usepackage{curves}
\usepackage{caption}
\usepackage{wasysym}
\usepackage{epsfig,epstopdf,graphicx}
\usepackage{tikz}
\usepackage{xcolor}
\usepackage{amsmath}
\usepackage{array}

\usetikzlibrary{arrows.meta, fit, backgrounds, positioning, calc, decorations.pathreplacing, shapes.geometric}

%% ── colour palette ──────────────────────────────────────────────────────────
\definecolor{active}{RGB}{52,152,219}       % blue  – active requests
\definecolor{waiting}{RGB}{231,76,60}       % red   – waiting requests
\definecolor{shared}{RGB}{39,174,96}        % green – shared-prefix chain
\definecolor{newblock}{RGB}{243,156,18}     % amber – newly-added block
\definecolor{evicted}{RGB}{192,57,43}       % dark red – evicted block
\definecolor{neutral}{RGB}{149,165,166}     % grey  – unrelated block
\definecolor{highlight}{RGB}{155,89,182}    % purple – highlighted result
\definecolor{bg}{RGB}{245,245,250}          % page bg

\setbeamercolor{background canvas}{bg=bg}
\setbeamertemplate{navigation symbols}{}
\setbeamerfont{title}{size=\large,series=\bfseries}

%% ── helpers ─────────────────────────────────────────────────────────────────
% Draw one hash block: \hashblock{x}{y}{label}{fillcolor}{textcolor}
\newcommand{\hashblock}[5]{%
  \node[draw, minimum width=1.1cm, minimum height=0.65cm,
        rounded corners=3pt, fill=#4, text=#5,
        font=\ttfamily\scriptsize, line width=0.6pt] at (#1,#2) {#3};
}

% Thick bracket above a range of blocks
% \chainbrace{x_start}{x_end}{y_top}{label}
\newcommand{\chainbrace}[4]{%
  \draw[decorate,
        decoration={brace,amplitude=5pt,raise=2pt},
        thick, color=shared]
    (#2+0.55,#3) -- (#1-0.55,#3) node[midway,above=7pt,color=shared,font=\small\bfseries]{#4};
}

\curvewarnfalse

\newtheorem{proposition}{Proposition}
\theoremstyle{example}
\newtheorem{theoremh}{Theorem}
\theoremstyle{plain}
\renewcommand{\textfraction}{0.01}
\renewcommand{\floatpagefraction}{0.99}
\newcommand{\ul}{\underline}
\newcounter{units}
\setbeamercovered{dynamic}
\logo{\includegraphics[width=0.5in,keepaspectratio]{\FIGDIR/iitb_logo.png}}
\newcommand{\FIGDIR}{../figures} 
\newcommand{\RESULTDIR}{../../results} 
\mode<presentation>{
    \usetheme[right,currentsection,hideothersubsections]{UTD}
    \useoutertheme{sidebar}
    \useinnertheme[shadow]{rounded}
    \usecolortheme{whale}
    \usecolortheme{orchid}
    \usefonttheme[onlymath]{serif}
    \setbeamertemplate{footline}{\centerline{Slide \insertframenumber/\inserttotalframenumber}}
}

\usebeamercolor[fg]{author in sidebar}
\title[{Systems for LLMs}]{\sc Benchmarking and Performance Analysis of LLM Inferencing Systems}
\author[\ul{Saksham Rathi}]{{\textbf{Saksham Rathi}}}
\institute[UTD]{\sc\small Systems for LLMs\\ Under Prof. Mythili Vutukuru}
\date[UCI]{Indian Institute of Technology Bombay \\ Autumn 2025-26}

\begin{document}
\frame{\titlepage}

\begin{frame}{Contents}
	\tableofcontents
\end{frame}

\section{The Core Trade-off}

\begin{frame}{The Core Trade-off}
    \begin{itemize}
        \item Larger Batches $\rightarrow$ Higher GPU Utilization
        \item Larger Shared Prefix $\rightarrow$ Better Temporal and Spatial Locality
        \item Need to ensure the prefix is the same across \textbf{all} the requests
    \end{itemize}
\end{frame}

\section{Chunked Hash Tree}

\begin{frame}{Changes to Chunked Hash Tree}
    \begin{itemize}
        \item Previously, we used to find the cost of adding a request to the working set (in terms of the newer chunks)
        \item Now we need information regarding the shared prefix chunks of all the active requests, along with changes to this upon adding a new request
        \item \textbf{Shared Prefix Chain:} A vector of hash blocks which is identical across all the active requests
        \item On activation: We iterate from back, and see if adding the request removes any blocks?
        \item On finishing: We iterate over the remaining blocks of the finished request, and check if any other blocks are shared across all the other active requests, and can be added to this chain. 
    \end{itemize}
\end{frame}



\section{RL Problem Formulation}

\begin{frame}{State}
    \begin{itemize}
        \item \textbf{B:} current batch size
        \item $\mathbf{C_B:}$ number of chunks shared before adding a new request
        \item $\mathbf{C_A:}$ number of chunks which will remain shared after adding the new (best) request
        \item $\mathbf{R_W:}$ number of requests waiting at the end of $C_A$
    \end{itemize}
\begin{block}{State}
    $s_t \in (B, C_B, C_A, R_W)$
\end{block}
\end{frame}

\begin{frame}{Action}
    \begin{itemize}
        \item \textbf{ADD:} pop and activate the best request from the chunked hash tree
        \item \textbf{SKIP:} break the batch construction loop, dispatch the batch to GPU
    \end{itemize}
\begin{block}{Action}
    $a_t \in (\texttt{ADD}, \texttt{SKIP})$
\end{block}
\end{frame}

\begin{frame}{Reward}
    \begin{itemize}
        \item Negative decode wall step time of the previous batch
        \item Delivered once per batch
    \end{itemize}
    \begin{block}{Reward}
    $r_t = -t_{\texttt{exec}}^{(t-1)}$
\end{block}
\end{frame}

\section{Policies}

% ------------------------------------------------
\begin{frame}{Policy 1: Q-Learning Scheduler}

    \begin{itemize}
        \item Trial-and-error learning method
        \item The system tries actions (ADD or SKIP), observes rewards and gradually learns which action works best in each situation.
        \item \textbf{Goal:} ``In this situation, which action gives the best long-term outcome?''
    \end{itemize}
\end{frame}
\begin{frame}{Policy 1: Q-Learning Scheduler}
\textbf{State Discretization (Making the problem manageable)}

The real system has continuous values.
We convert them into buckets (like "low", "medium", "high").

\[
\tilde{s}_t = (\tilde{B}, \widetilde{C_B - C_A}, \tilde{R}_W)
\]

Where:
\begin{itemize}
    \item $B$ = batch size
    \item $C_B - C_A$ = chunk difference (loss signal)
    \item $R_W$ = number of waiting requests
\end{itemize}

\textbf{Intuition:}
Instead of remembering every exact number,
we group similar situations together.

\end{frame}
\begin{frame}{Policy 1: Q-Learning Scheduler}

\textbf{Q-Function (What is being learned?)}

\[
Q(s_t, a_t)
\]

This is a table storing:
\[
Q : (\tilde{s}, a) \rightarrow \mathbb{R}
\]

For each state and action,
store a number that estimates
``how good this action is in this situation''.

\end{frame}

% ------------------------------------------------
\begin{frame}{Policy 1: Q-Learning Scheduler}

The system must balance:
\begin{itemize}
    \item Try new actions (exploration)
    \item Use what it already learned (exploitation)
\end{itemize}

\vspace{0.3cm}

\textbf{Epsilon-Greedy Policy}

\[
a_t =
\begin{cases}
\text{random action} & \text{with probability } \epsilon \\
\arg\max_{a} Q(s_t, a) & \text{otherwise}
\end{cases}
\]

\begin{itemize}
    \item With probability $\epsilon$, try something random
    \item Otherwise, choose the action with highest learned value
\end{itemize}

\end{frame}

% ------------------------------------------------
\begin{frame}{Policy 1: Q-Learning Scheduler}
\textbf{Intuition:}
Early on, explore a lot.
Later, trust what you learned.


\textbf{Exploration Decay}

\[
\epsilon_t =
\begin{cases}
\epsilon_{\text{start}} -
(\epsilon_{\text{start}} - \epsilon_{\text{end}})
\frac{t}{T_{\text{decay}}}
& t < T_{\text{decay}} \\
\epsilon_{\text{end}} & \text{otherwise}
\end{cases}
\]


\end{frame}

% ------------------------------------------------
\begin{frame}{Policy 1: Q-Learning Scheduler}

\textbf{How Learning Happens}

After taking action $a_t$:
\begin{itemize}
    \item Observe reward $r_t$
    \item Update the table
\end{itemize}


\textbf{Update Rule}

\[
Q(s_t, a_t)
\leftarrow
Q(s_t, a_t)
+
\alpha
\left(
r_t - Q(s_t, a_t)
\right)
\]

\begin{itemize}
    \item If reward is better than expected → increase value
    \item If worse → decrease value
\end{itemize}

$\alpha$ = learning rate (how fast we adjust)

\end{frame}

% ------------------------------------------------
\begin{frame}{Policy 1: Q-Learning Scheduler}

\textbf{Overall Objective}

\[
\max_\pi \; \mathbb{E}
\left[
\sum_{t=0}^{\infty}
\gamma^t r_t
\right]
\]

Find a policy $\pi$ that maximizes long-term rewards.

Not just immediate reward, we care about the future too.

\end{frame}

% ------------------------------------------------
\begin{frame}{Policy 2: Contextual Bandit (UCB)}

    \begin{itemize}
        \item This is a simpler learning strategy.
        \item It does NOT consider long-term future.
        \item \textbf{Goal:} ``Which action works best right now in this situation?''
    \end{itemize}

\[
\hat{s}_t =
\text{bin}(B, C_B - C_A, R_W)
\]

We still bucket the state.

\end{frame}

% ------------------------------------------------
\begin{frame}{Policy 2: Contextual Bandit (UCB)}

\textbf{Statistics Stored}

\[
N(\hat{s}, a) \quad \text{(how many times tried)}
\]

\[
\bar{r}(\hat{s}, a) =
\frac{\sum r}{N(\hat{s}, a)}
\]

Track average reward of each action.

\end{frame}

% ------------------------------------------------
\begin{frame}{Policy 2: Contextual Bandit (UCB)}

\textbf{Upper Confidence Bound (UCB)}

\[
\text{UCB}(\hat{s}, a)
=
\bar{r}(\hat{s}, a)
+
c
\sqrt{
\frac{\log T}{N(\hat{s}, a)}
}
\]

\textbf{Intuition:}

Two parts:
\begin{itemize}
    \item Exploit: choose high average reward
    \item Explore: boost rarely tried actions
\end{itemize}

If an action was rarely tried,
its uncertainty is large → encourage exploration.

\end{frame}

% ------------------------------------------------
\begin{frame}{Policy 2: Contextual Bandit (UCB)}

\textbf{Policy}

\[
a_t = \arg\max_a \text{UCB}(\hat{s}_t, a)
\]

\textbf{Meaning:}
Choose the action with highest optimistic estimate.

\end{frame}

% ------------------------------------------------
\begin{frame}{Policy 3: Heuristic Scheduler}

    \begin{itemize}
        \item This policy does NOT learn.
        \item It uses manually designed rules.
    \end{itemize}

\vspace{1cm}

\textbf{Chunk Loss} $\Delta C = C_B - C_A$

Interpretation:
\begin{itemize}
    \item If $\Delta C$ small $\rightarrow$ memory accesses remain similar
    \item If large $\rightarrow$ too much loss in temporal and spatial locality
\end{itemize}


\end{frame}

% ------------------------------------------------
\begin{frame}{Policy 3: Heuristic Scheduler}

\textbf{Decision Logic}

\begin{enumerate}
    \item If no batch $\rightarrow$ ADD
    \item If no chunk loss $\rightarrow$ ADD
    \item If batch small $\rightarrow$ tolerate slightly more loss
    \item If batch large $\rightarrow$ be stricter
    \item If many waiting requests $\rightarrow$ prioritize ADD
    \item Otherwise $\rightarrow$ SKIP
\end{enumerate}
\end{frame}


\section{}
\begin{frame}
    \Huge{\centerline{\bf Thank You}}
\end{frame}

\end{document}